\documentclass[code]{mancls}
\title{SUSTech-fiBiD 中国证券市场全尺度数字“风洞” V1.0}
\usepackage{listings}
\lstset{language=Python}

\begin{document}

\section{说明}
本源程序文档采用标准格式收录本项目的核心源代码文件，覆盖仿真内核、订单与撮合、消息与时钟、交易所回测接口、智能体框架、自校准与指标等关键模块。为便于审阅，统一使用 \texttt{codestyleln} 风格显示行号。

\section{核心仿真内核}
\subsection{内核调度与事件循环（core/kernel.py）}
\lstinputlisting[style=codestyleln]{../../core/kernel.py}

\subsection{时钟与时间推进（core/clock.py）}
\lstinputlisting[style=codestyleln]{../../core/clock.py}

\subsection{消息结构（core/message.py）}
\lstinputlisting[style=codestyleln]{../../core/message.py}

\section{订单、订单簿与撮合}
\subsection{订单结构定义（core/order.py）}
\lstinputlisting[style=codestyleln]{../../core/order.py}

\section{股票模块实现}
\subsection{股票数据管理（core/symbol.py）}
\lstinputlisting[style=codestyleln]{../../core/symbol.py}

\subsection{订单簿实现（core/orderbook.py）}
\lstinputlisting[style=codestyleln]{../../core/orderbook.py}

\subsection{轻量 LOB 辅助（core/lob.py）}
\lstinputlisting[style=codestyleln]{../../core/lob.py}

\section{智能体框架与示例}
\subsection{代理基类（core/agent/base.py）}
\lstinputlisting[style=codestyleln]{../../core/agent/base.py}

\subsection{零智商代理（core/agent/zero\_intelligence.py）}
\lstinputlisting[style=codestyleln]{../../core/agent/zero_intelligence.py}

\subsection{噪声代理（core/agent/noise.py）}
\lstinputlisting[style=codestyleln]{../../core/agent/noise.py}

\section{自校准与指标}
\subsection{校准上下文（core/calibration/context.py）}
\lstinputlisting[style=codestyleln]{../../core/calibration/context.py}

\subsection{校准指标（core/calibration/metrics.py）}
\lstinputlisting[style=codestyleln]{../../core/calibration/metrics.py}

\section{日志与工具}
\subsection{日志模块（core/logger.py）}
\lstinputlisting[style=codestyleln]{../../core/logger.py}

\subsection{通用工具（util/util.py）}
\lstinputlisting[style=codestyleln]{../../util/util.py}

\end{document}
